% !TEX program = xelatex
\documentclass[12pt,a4paper]{article}

\usepackage{xeCJK}
\setCJKmainfont{Hiragino Mincho ProN}
\setCJKsansfont{Hiragino Kaku Gothic ProN}

\usepackage{amsmath,amssymb}
\usepackage{geometry}
\geometry{margin=2.5cm}
\usepackage{hyperref}
\hypersetup{colorlinks=true,linkcolor=blue,citecolor=blue,urlcolor=blue}
\usepackage{booktabs}
\usepackage{array}
\usepackage{tcolorbox}
\tcbuselibrary{skins,breakable}
\usepackage{graphicx}
\usepackage{enumitem}

\newtcolorbox{keypoint}[1][]{
  colback=blue!5!white,
  colframe=blue!50!black,
  fonttitle=\bfseries,
  title={#1},
  breakable
}

\newtcolorbox{funfact}[1][]{
  colback=green!5!white,
  colframe=green!50!black,
  fonttitle=\bfseries,
  title={#1},
  breakable
}

\newtcolorbox{honesty}[1][]{
  colback=red!5!white,
  colframe=red!50!black,
  fonttitle=\bfseries,
  title={#1},
  breakable
}

\newtcolorbox{analogy}[1][]{
  colback=orange!5!white,
  colframe=orange!50!black,
  fonttitle=\bfseries,
  title={#1},
  breakable
}

\title{\textbf{なぜ宇宙は4次元なのか？}\\
{\Large 整数の世界から読み解く時空の秘密}\\[0.5em]
{\normalsize --- 高校生のためのガイド ---}}
\author{ライト兄弟コラボレーション}
\date{2026年4月}

\begin{document}
\maketitle

\begin{keypoint}[この文書について]
この文書は、「なぜ宇宙は空間3次元＋時間1次元＝4次元なのか？」という
根本的な疑問に対して、ライト兄弟 (Wright Brothers) の枠組みがどのように
答えるかを、高校生向けにわかりやすく説明するものです。
数学の予備知識として、素数・対数・複素数の基礎があれば十分です。
\end{keypoint}

\tableofcontents
\newpage

%% ============================================================
\section{問い：なぜ私たちは「3＋1＝4次元」に住んでいるのか？}
%% ============================================================

\subsection{当たり前のようで、当たり前でないこと}

私たちは日常的に3つの空間方向（前後・左右・上下）と1つの時間方向（過去から未来）を
経験しています。合計4次元。これは「当たり前」のように思えますが、
実は物理学最大の謎の一つです。

なぜ5次元でも6次元でもなく、ちょうど4次元なのでしょうか？

\begin{analogy}[家のアナロジー]
家を設計するとき、「3部屋＋1つの廊下」という構造を考えてみましょう。
\begin{itemize}
  \item \textbf{3部屋}（空間）：自由に動き回れる場所。前後・左右・上下。
  \item \textbf{1つの廊下}（時間）：一方通行で、戻れない通路。過去から未来へ。
\end{itemize}
なぜ5部屋でも2部屋でもなく、ちょうど3部屋なのか？
なぜ廊下は1本なのか？ これが宇宙版の設計問題です。
\end{analogy}

\subsection{Ehrenfestの観察（1917年）}

オランダの物理学者Ehrenfest（エーレンフェスト）は1917年に、
次元の数が物理法則にどう影響するかを調べました。

\begin{keypoint}[Ehrenfestの発見]
\begin{itemize}
  \item \textbf{空間が4次元以上}だと、惑星の軌道は不安定になる。
    地球は太陽の周りを回り続けることができず、太陽に落ちるか飛び去ってしまう。
  \item \textbf{空間が2次元以下}だと、原子が安定に存在できない。
    電子が原子核を周回するための十分な「部屋」がない。
  \item \textbf{空間がちょうど3次元}のときだけ、
    安定な軌道と安定な原子の両方が可能。
\end{itemize}
\end{keypoint}

Ehrenfestは「3次元でないと私たちは存在できない」ことを示しました。
しかし、「なぜ宇宙が3次元を\textbf{選んだ}のか」は答えていません。
「人間が住めるから3次元」（人間原理）ではなく、
「3次元でなければ\textbf{ならない}数学的理由」があるのでしょうか？

ライト兄弟の答え：\textbf{あります。整数の構造がそれを決めています。}

%% ============================================================
\section{$\mathrm{Spec}(\mathbb{Z})$ とは何か？ --- 整数のDNA}
%% ============================================================

\subsection{素数：整数の建築ブロック}

すべての正整数は、素数の積として一意に表せます（算術の基本定理）：
\[
  12 = 2^2 \times 3, \quad 30 = 2 \times 3 \times 5, \quad 2026 = 2 \times 1013.
\]

素数は整数の「原子」です。化学で物質が原子からできているように、
数は素数からできています。

\subsection{Spec(Z)：素数の世界地図}

数学者は「$\mathrm{Spec}(\mathbb{Z})$（スペック・ゼット）」という対象を
考えます。これは整数の環 $\mathbb{Z}$ の「スペクトル」で、
大雑把に言えば：

\begin{keypoint}[$\mathrm{Spec}(\mathbb{Z})$ のイメージ]
$\mathrm{Spec}(\mathbb{Z})$ とは、素数 $2, 3, 5, 7, 11, \ldots$ を
「点」として持つ空間です。さらに「$(0)$」という特別な点
（有理数 $\mathbb{Q}$ に対応）もあります。
\end{keypoint}

\begin{analogy}[鉄道路線図のアナロジー]
$\mathrm{Spec}(\mathbb{Z})$ を鉄道路線図だと思ってください。
\begin{itemize}
  \item 各\textbf{駅}が素数：2番駅、3番駅、5番駅、7番駅、\ldots
  \item 駅同士を結ぶ\textbf{路線}が、素数間の算術的関係。
  \item 路線図全体（$\mathrm{Spec}(\mathbb{Z})$）が、整数の構造を表す「空間」。
\end{itemize}
この路線図は紙の上（2次元）に描けますが、数学的な意味での「厚み」は
3次元であることが1960年代に証明されました。これが次の話題です。
\end{analogy}

\subsection{なぜ $\mathrm{Spec}(\mathbb{Z})$ が重要なのか}

ライト兄弟の枠組みでは、こう主張します：

\begin{keypoint}[基本仮定]
宇宙の空間構造は、$\mathrm{Spec}(\mathbb{Z})$ の幾何学に基づいている。
\end{keypoint}

つまり、空間の次元数は $\mathrm{Spec}(\mathbb{Z})$ の「数学的な次元」で
決まるということです。これは大胆な仮定ですが、驚くべき結果を生みます。

%% ============================================================
\section{$\mathrm{Spec}(\mathbb{Z})$ は「3次元的」（Artin--Verdierの定理）}
%% ============================================================

\subsection{コホモロジー次元：空間の「厚み」を測る}

1960年代に、数学者のArtin（アルティン）とVerdier（ヴェルディエ）は
驚くべき定理を証明しました。

\begin{keypoint}[Artin--Verdierの定理（1964年）]
$\mathrm{Spec}(\mathbb{Z})$ の「エタールコホモロジー次元」はちょうど\textbf{3}である。
\[
  \mathrm{cd}_{\text{\'et}}(\mathrm{Spec}(\mathbb{Z})) = 3
\]
\end{keypoint}

「エタールコホモロジー次元」は難しそうな言葉ですが、
直感的には「その空間が何次元の図形のように振る舞うか」を測る数です。

\begin{analogy}[影のアナロジー]
物体の「次元」を知りたいとき、いろいろな角度から光を当てて影を調べる方法があります。
\begin{itemize}
  \item 1次元の物体（線）：影は点か線。
  \item 2次元の物体（面）：影は線か面。
  \item 3次元の物体（立体）：影は面。
\end{itemize}
コホモロジーは、この「影の調べ方」の数学的に厳密なバージョンです。
Artin--Verdierの定理は、$\mathrm{Spec}(\mathbb{Z})$ に数学的な光を当てると、
\textbf{3次元の物体}のように影が落ちることを示しています。
\end{analogy}

重要な点：これは\textbf{予想}ではなく、60年以上前に\textbf{証明された定理}です。
数学の世界では完全に確立された結果です。

\subsection{算術的位相幾何学：素数は結び目}

さらに面白いことに、数学者のMazur（メイザー）やMorishita（森下昌紀）は、
$\mathrm{Spec}(\mathbb{Z})$ と3次元多様体の間に深いアナロジーがあることを
発見しました。

\begin{funfact}[素数＝結び目]
\begin{center}
\begin{tabular}{ll}
\toprule
\textbf{整数の世界} & \textbf{3次元空間の世界} \\
\midrule
$\mathrm{Spec}(\mathbb{Z})$ & 3次元の空間 $M^3$ \\
素数 $p$ & 空間内の結び目（ロープの輪） \\
素数間の関係 & 結び目の絡まり方 \\
\bottomrule
\end{tabular}
\end{center}
素数は、3次元空間の中にある「結び目」のようなものなのです！
\end{funfact}

%% ============================================================
\section{なぜ「3」なのか：$3 = 1 + 2$ の分解}
%% ============================================================

「コホモロジー次元が3」と言われても、なぜちょうど3なのかが気になります。
実は、この「3」は2つの部分に分解できます。

\subsection{第1の寄与：整数の「基線」= 1}

整数 $\mathbb{Z}$ の中の素イデアル（素数に対応する特別な構造）は、
次のような「鎖」を作ります：
\[
  (0) \subset (p)
\]
つまり、「一般的な場所 $(0)$」と「特定の素数の場所 $(p)$」の2段階。
この鎖の長さが1なので、$\mathbb{Z}$ のKrull次元は\textbf{1}です。

\begin{analogy}[数直線のアナロジー]
整数を数直線上に並べると1次元。これがKrull次元 = 1の直感的な意味です。
整数の基本的な構造は「1本の線」のようなものなのです。
\end{analogy}

\subsection{第2の寄与：$\sqrt{-1}$ の存在 = 2}

ここからが面白いところです。実数 $\mathbb{R}$ を超えて複素数 $\mathbb{C}$ に
進むと、$\sqrt{-1} = i$ という新しい数が登場します。

$\mathbb{C}$ から $\mathbb{R}$ に戻る操作は「複素共役」です：
$a + bi \mapsto a - bi$。この操作は $\mathbb{C}$ の自己同型群（対称性）を作り、
これをガロア群と呼びます：
\[
  \mathrm{Gal}(\mathbb{C}/\mathbb{R}) = \mathbb{Z}/2
\]

この $\mathbb{Z}/2$（2つの元からなる群：「何もしない」と「複素共役する」）が、
コホモロジー次元に\textbf{2}を寄与します。

\begin{keypoint}[3の分解]
\[
  3 = \underbrace{1}_{\text{Krull次元（整数の基線）}} +
      \underbrace{2}_{\text{$\mathrm{Gal}(\mathbb{C}/\mathbb{R}) = \mathbb{Z}/2$ の寄与}}
\]
\begin{itemize}
  \item \textbf{1}：整数が「1本の線」であることから。
  \item \textbf{2}：$\sqrt{-1}$ が存在する（複素数が実数の2次拡大である）ことから。
\end{itemize}
空間が3次元なのは、整数の基線（1次元）に、
$\sqrt{-1}$ による2次元が加わるからです！
\end{keypoint}

そして注目してください：ここに現れた数「2」は、ガロア群 $\mathbb{Z}/2$ の
「2」です。この「2」がこの後、何度も何度も登場します。

%% ============================================================
\section{時間の追加：一方通行のドア}
%% ============================================================

\subsection{空間と時間の違い}

空間では前後左右上下、自由に動けます。しかし時間は違います：
\begin{itemize}
  \item 過去には戻れない（時間の矢）。
  \item 始まりがある（ビッグバン）。
  \item 未来は開かれている。
\end{itemize}

この非対称性が、時間が空間と質的に異なる理由です。

\subsection{Wheeler--DeWitt方程式：宇宙のシュレーディンガー方程式}

量子力学では、粒子の状態をシュレーディンガー方程式で記述します。
宇宙全体にも同様の方程式があります。それがWheeler--DeWitt (WDW) 方程式です：
\[
  \hat{H}\,\Psi = 0
\]
ここで $\Psi$ は「宇宙の波動関数」、$\hat{H}$ は宇宙のハミルトニアンです。

\subsection{ディリクレ・ノイマン境界条件}

WDW方程式に以下の境界条件を課します：

\begin{keypoint}[DN境界条件]
\begin{itemize}
  \item \textbf{ディリクレ条件}（始まり）：宇宙の始まり（ビッグバン）で
    波動関数はゼロ。$\Psi|_{a=0} = 0$。

    イメージ：「壁にしっかり固定されたロープの端」
  \item \textbf{ノイマン条件}（終わり）：宇宙の未来は開かれている。
    $\partial_a \Psi|_{a \to \infty} = 0$。

    イメージ：「自由に揺れるロープの端」
\end{itemize}
\end{keypoint}

\begin{analogy}[ロープのアナロジー]
ギターの弦を想像してください。
\begin{itemize}
  \item 片方の端は固定（ディリクレ）：ビッグバン。動けない。
  \item もう片方の端は自由（ノイマン）：未来。自由に振動できる。
\end{itemize}
この非対称性が「時間」を生み出します。
固定端→自由端の方向が「時間の流れ」に対応します。
\end{analogy}

\subsection{結果：3＋1＝4}

\begin{keypoint}[時空次元の公式]
\[
  d_{\text{時空}} = \underbrace{3}_{\mathrm{Spec}(\mathbb{Z})\text{のコホモロジー次元}}
  + \underbrace{1}_{\text{WDW DN境界条件}} = 4
\]
\begin{itemize}
  \item 空間の3次元：整数の算術構造（Artin--Verdierの定理）から。
  \item 時間の1次元：宇宙に「始まり」と「開かれた未来」があることから。
\end{itemize}
\end{keypoint}

%% ============================================================
\section{すべては $\mathbb{Z}/2$（数「2」）に帰着する}
%% ============================================================

ここまでの議論を振り返ると、驚くべきことに気づきます。
あちこちに「2」が現れているのです。

\begin{keypoint}[「2」の遍在性]
\begin{center}
\renewcommand{\arraystretch}{1.3}
\begin{tabular}{p{6cm}l}
\toprule
\textbf{どこに現れるか} & \textbf{「2」の役割} \\
\midrule
$\mathrm{Gal}(\mathbb{C}/\mathbb{R}) = \mathbb{Z}/2$ &
空間の「3」に2を寄与（$3=1+2$）\\
\addlinespace
$p = 2$（最小の素数）&
ダークエネルギー $\Omega_\Lambda = 2\pi/9$ を決定 \\
\addlinespace
$\sqrt{-1}$ の存在 &
$\mathbb{C}$ は $\mathbb{R}$ の\textbf{2}次拡大 \\
\addlinespace
スピン &
粒子は $\mathbb{Z}/2$ 的性質（上向き/下向き）を持つ \\
\addlinespace
$2^4 = 4^2$ &
2の「べき乗双子」が次元4を与える \\
\bottomrule
\end{tabular}
\end{center}

すべての道は「2」に通じます。
\end{keypoint}

なぜ「2」がこれほど特別なのでしょうか？ それは「2」が最小の素数だからです。
素数の中で最も単純で基本的な数。宇宙の構造を決める「種」のような存在です。

%% ============================================================
\section{魔法の方程式：$2^4 = 4^2$}
%% ============================================================

\subsection{$p^n = n^p$ の非自明解}

次の方程式を考えましょう：
\[
  p^n = n^p \quad (p \neq n, \text{ ともに正整数})
\]

これは「$p$ の $n$ 乗と $n$ の $p$ 乗が等しい」という条件です。
いくつか試してみましょう：

\begin{center}
\renewcommand{\arraystretch}{1.2}
\begin{tabular}{ccccl}
\toprule
$p$ & $n$ & $p^n$ & $n^p$ & 等しい？ \\
\midrule
2 & 3 & 8 & 9 & 惜しい！でも不等 \\
2 & 4 & \textbf{16} & \textbf{16} & \textbf{等しい！} \\
2 & 5 & 32 & 25 & 不等 \\
3 & 4 & 81 & 64 & 不等 \\
3 & 5 & 243 & 125 & 不等 \\
3 & 27 & $3^{27}$ & $27^3 = 19683$ & 全く違う \\
\bottomrule
\end{tabular}
\end{center}

\begin{keypoint}[定理]
$p \neq n$ なる正整数の中で、$p^n = n^p$ が成り立つのは
$\{p, n\} = \{2, 4\}$の場合——つまり $2^4 = 4^2 = 16$ の場合——
\textbf{だけ}です。
\end{keypoint}

\subsection{証明のアイデア（高校数学で理解できます！）}

関数 $f(x) = \frac{\ln x}{x}$ を考えます。$p^n = n^p$ を対数で書くと：
\[
  n \ln p = p \ln n \quad \Longleftrightarrow \quad \frac{\ln p}{p} = \frac{\ln n}{n}
  \quad \Longleftrightarrow \quad f(p) = f(n)
\]

つまり、$f(x)$ のグラフで同じ高さにある2つの異なる正整数を見つける問題です。

$f(x)$ の微分は：
\[
  f'(x) = \frac{1 - \ln x}{x^2}
\]
\begin{itemize}
  \item $x < e \approx 2.718$ で $f'(x) > 0$（増加）
  \item $x > e$ で $f'(x) < 0$（減少）
  \item $x = e$ で最大値
\end{itemize}

$f(p) = f(n)$（$p \neq n$）が成り立つには、一方が $e$ より左、
もう一方が $e$ より右にある必要があります。
$e$ より左にある正整数は $1$ と $2$ だけ。

$p = 1$ の場合は自明解（$1^n = n^1$ なら $n=1$）。

$p = 2$ の場合、$f(2) = \frac{\ln 2}{2}$ と等しくなる $n > e$ を探すと、
$n = 4$ がぴったり合います：
\[
  f(4) = \frac{\ln 4}{4} = \frac{2\ln 2}{4} = \frac{\ln 2}{2} = f(2) \quad \checkmark
\]

\begin{funfact}[べき乗双子]
「2」は唯一の「べき乗双子」を持つ正整数です。その双子は「4」。
\[
  2^{\mathbf{4}} = \mathbf{4}^2 = 16
\]
他のどの数にも、こんなパートナーはいません。
そして $4 = $ 時空の次元数なのです！
\end{funfact}

%% ============================================================
\section{四元数：4次元の数体系}
%% ============================================================

\subsection{数の体系の階段}

数学には「ノルム付き分割代数」と呼ばれる特別な数体系があります。
1898年にHurwitz（フルヴィッツ）は、そのような体系は\textbf{4種類しかない}ことを
証明しました。

\begin{center}
\renewcommand{\arraystretch}{1.3}
\begin{tabular}{lccc}
\toprule
\textbf{名前} & \textbf{次元} & \textbf{かけ算の順序} & \textbf{かけ算の結合法則} \\
\midrule
実数 $\mathbb{R}$ & $1 = 2^0$ & 交換可能 & 結合的 \\
複素数 $\mathbb{C}$ & $2 = 2^1$ & 交換可能 & 結合的 \\
四元数 $\mathbb{H}$ & $4 = 2^2$ & \textbf{交換不可} & 結合的 \\
八元数 $\mathbb{O}$ & $8 = 2^3$ & 交換不可 & \textbf{結合的でない} \\
\bottomrule
\end{tabular}
\end{center}

\begin{analogy}[数の階段]
\begin{itemize}
  \item \textbf{実数}（1次元）：数直線。何でもできる優等生。
  \item \textbf{複素数}（2次元）：平面。$a + bi$。まだ $ab = ba$ が成り立つ。
  \item \textbf{四元数}（4次元）：$a + bi + cj + dk$。ここで $ij \neq ji$
    （かけ算の順序が大事になる）。でも$(ab)c = a(bc)$はまだ成り立つ。
  \item \textbf{八元数}（8次元）：かけ算の順序も結合法則も崩壊。数学的に「野性的」。
\end{itemize}
階段を上るたびに、何かが失われていきます。
\end{analogy}

\subsection{四元数の特別さ}

四元数 $\mathbb{H}$ は「最後の行儀の良い数体系」です。

\begin{keypoint}[四元数と時空]
\begin{itemize}
  \item 四元数の次元 = $4$ = 時空の次元
  \item 四元数の次元 = $p^2 = 2^2 = 4$（特別な素数 $p=2$ の平方）
  \item 四元数は3次元の回転を記述する（$\mathrm{SU}(2) \cong$ 単位四元数）
  \item 3D回転 → 空間3次元と四元数4次元が自然に結びつく
\end{itemize}
時空の次元と「最も良い」数体系の次元が一致するのは偶然でしょうか？
\end{keypoint}

ゲームやCGで3Dの回転を扱うとき、実はプログラマーは四元数を使っています。
3次元空間の回転を最も効率的に記述できるのが、4次元の数体系（四元数）なのです。
これは空間3次元と数体系4次元の間の深い関連を示唆しています。

%% ============================================================
\section{ダークエネルギーとの関連}
%% ============================================================

\subsection{ダークエネルギーとは}

宇宙の約70\%は「ダークエネルギー」と呼ばれる謎のエネルギーで満たされています。
このエネルギーが宇宙の膨張を加速させています。

\subsection{同じ「2」が答える}

ライト兄弟の枠組みでは、ダークエネルギーの量も計算できます：
\[
  \Omega_\Lambda = \frac{2\pi}{9} \approx 0.698
\]

観測値は $\Omega_\Lambda \approx 0.689$。非常によく合っています！

\begin{keypoint}[統一的な答え]
\begin{center}
\begin{tabular}{ll}
\toprule
\textbf{問い} & \textbf{答えの源} \\
\midrule
なぜ空間は3次元？ & $\mathrm{Gal}(\mathbb{C}/\mathbb{R}) = \mathbb{Z}/2$
→ $\mathrm{cd} = 1+2 = 3$ \\
なぜ時空は4次元？ & $3 + 1 = 4$（WDW DN境界条件で $+1$）\\
なぜダークエネルギーはこの値？ & 同じ $p=2$ → $\Omega_\Lambda = 2\pi/9$ \\
\bottomrule
\end{tabular}
\end{center}

「なぜ4次元？」と「なぜこのダークエネルギー？」は別々の問いのように見えて、
\textbf{実は同じ答え}（$\mathbb{Z}/2$ と $p=2$）から出ているのです。
\end{keypoint}

これは非常に美しい統一です。2つの全く異なる物理現象——時空の次元数と
宇宙のダークエネルギー——が、同じ数学的源泉から流れ出ているのです。

%% ============================================================
\section{弦理論との比較}
%% ============================================================

\subsection{弦理論のアプローチ}

弦理論は現代物理学の最も有力な理論の一つです。しかし、次元に関しては：

\begin{itemize}
  \item 弦理論は宇宙が\textbf{10次元}（または11次元、26次元）であると予言する。
  \item 私たちが見えるのは4次元。残りの6次元は？
  \item 答え：「余分な6次元はCalabi--Yau多様体という微小な形に丸まっている」
\end{itemize}

\subsection{折り紙 vs 設計図}

\begin{analogy}[折り紙 vs 設計図]
\textbf{弦理論}は「折り紙」アプローチ：
\begin{itemize}
  \item 10次元の大きな紙（時空）を用意する
  \item 6次元分を複雑に折りたたむ（コンパクト化）
  \item 残った4次元が私たちの宇宙
  \item 問題：折り方が $10^{500}$ 通りもある！（弦理論のランドスケープ問題）
\end{itemize}

\textbf{ライト兄弟}は「設計図」アプローチ：
\begin{itemize}
  \item 整数 $\mathbb{Z}$ の設計図（$\mathrm{Spec}(\mathbb{Z})$）を読む
  \item 設計図は「3次元の空間＋1次元の時間」と書いてある
  \item 最初から4次元。折りたたむ必要なし
  \item 余分な次元も、膨大な選択肢もない
\end{itemize}
\end{analogy}

\begin{center}
\renewcommand{\arraystretch}{1.3}
\begin{tabular}{p{5.5cm}p{5.5cm}}
\toprule
\textbf{弦理論} & \textbf{ライト兄弟} \\
\midrule
基本次元 = 10（または11、26） & 基本次元 = 4 \\
余分な次元をコンパクト化 & コンパクト化不要 \\
$10^{500}$ 通りの可能性 & 一意な結果 \\
「なぜ4次元？」→ 人間原理 & 「なぜ4次元？」→ 算術的必然 \\
\bottomrule
\end{tabular}
\end{center}

\begin{funfact}[公平な注意]
弦理論は素粒子物理学で多くの成果を上げており、
ライト兄弟の枠組みとは矛盾するのではなく、異なる視点を提供するものです。
どちらが正しいかは、最終的には実験と観測が判断します。
\end{funfact}

%% ============================================================
\section{正直な限界}
%% ============================================================

科学で大切なのは、自分の理論の限界を正直に認めることです。

\begin{honesty}[認めるべき限界]
\begin{enumerate}[label=\textbf{\arabic*.}]
  \item \textbf{「空間 $= \mathrm{Spec}(\mathbb{Z})$」は仮定であり、証明ではない。}

  なぜ空間が整数のスペクトルに基づくのかは、この理論の出発点（公理）であり、
  より深い原理から導かれたものではありません。
  これは「なぜ？」の問いをもう一段深く先送りしているとも言えます。

  \item \textbf{なぜ $3+1$ であり、$2+2$ ではないのか？}

  時空が4次元であることは導けますが、なぜ空間3＋時間1であり、
  空間2＋時間2（超双曲的符号）ではないのかは完全には説明できていません。
  DN境界条件が時間を空間と区別する手がかりを与えますが、
  厳密な導出はまだ未完成です。

  \item \textbf{「時間に $+1$」はやや場当たり的}

  空間の3次元は美しく算術から導かれますが、
  時間の1次元の追加はWDW方程式の境界条件という別の議論に依存しており、
  同じ算術構造から統一的に出てくるわけではありません。

  \item \textbf{これは枠組みであり、最終理論ではない}

  ライト兄弟は「整数の算術から物理を導く」という研究プログラムです。
  多くの予測は成功していますが、標準模型のすべてを導出できているわけではなく、
  まだ発展途上の枠組みです。
\end{enumerate}
\end{honesty}

これらの限界は、研究の終わりではなく、\textbf{次の研究の始まり}です。
科学は完璧な答えを出すことではなく、よりよい問いを立て続けることなのです。

%% ============================================================
\section{メッセージ：宇宙の次元は整数が決めている}
%% ============================================================

最後に、この章全体のメッセージをまとめましょう。

\begin{keypoint}[まとめ]
\textbf{なぜ宇宙は4次元なのか？}

\begin{enumerate}[label=\textbf{Step \arabic*:}]
  \item 整数の世界 $\mathrm{Spec}(\mathbb{Z})$ は、数学的に「3次元」である
    （Artin--Verdierの定理、1964年に証明済み）。
  \item この「3」は $3 = 1 + 2$ と分解される。
    $1$ は整数の基線、$2$ は $\sqrt{-1}$ の存在（$\mathrm{Gal}(\mathbb{C}/\mathbb{R}) = \mathbb{Z}/2$）から。
  \item 宇宙には始まり（ビッグバン）があり未来は開かれている。
    この非対称性が1つの時間次元を追加する。
  \item 結果：$d = 3 + 1 = 4$。
\end{enumerate}

すべては $\mathbb{Z}/2$——最小の素数「2」の群——に帰着します。
同じ「2」がダークエネルギーの値 $\Omega_\Lambda = 2\pi/9$ も決めています。
\end{keypoint}

\begin{funfact}[最後に]
\textbf{宇宙の次元は、「2が最小の素数である」という事実と同じくらい不可避
かもしれません。}

素数の中で最も小さく、最も特別な数「2」。この数が：
\begin{itemize}
  \item 空間を3次元にし（$\mathbb{Z}/2$ が2を寄与）
  \item 時空を4次元にし（$3+1=4$）
  \item ダークエネルギーの値を決め（$p=2$ → $\Omega_\Lambda = 2\pi/9$）
  \item 自分の「べき乗双子」として4を生み出す（$2^4 = 4^2$）
\end{itemize}

私たちが住む宇宙は、整数のDNAに書き込まれた必然なのかもしれません。

次にテストで $2^4 = 16 = 4^2$ を見たとき、そこに宇宙の設計図を見てください。
\end{funfact}

\vspace{2em}

\begin{center}
\Large\textbf{--- おわり ---}
\end{center}

\vspace{1em}

\subsection*{さらに学びたい人へ}

\begin{itemize}
  \item \textbf{素数について}：
    Marcus du Sautoy『素数の音楽』（新潮文庫）——素数の不思議さを一般向けに解説。
  \item \textbf{次元について}：
    Ian Stewart『次元についてのお話』——なぜ空間は3次元かを探求。
  \item \textbf{複素数について}：
    Paul Nahin『An Imaginary Tale: The Story of $\sqrt{-1}$』——$\sqrt{-1}$ の歴史。
  \item \textbf{四元数について}：
    ゲーム開発やCG分野の教科書に四元数（クォータニオン）の実用的な解説がある。
  \item \textbf{算術的位相幾何学について}：
    森下昌紀『結び目と素数』——素数と結び目のアナロジーを解説する専門書。
\end{itemize}

\subsection*{用語集}

\begin{center}
\renewcommand{\arraystretch}{1.3}
\begin{tabular}{p{4cm}p{8cm}}
\toprule
\textbf{用語} & \textbf{意味} \\
\midrule
$\mathrm{Spec}(\mathbb{Z})$ & 整数環のスペクトル。素数を「点」とする空間。 \\
コホモロジー次元 & 空間の「数学的な厚み」を測る数。 \\
Artin--Verdier双対性 & $\mathrm{Spec}(\mathbb{Z})$ のコホモロジー次元が3であることを示す定理。 \\
$\mathrm{Gal}(\mathbb{C}/\mathbb{R})$ & 複素数から実数へのガロア群。$= \mathbb{Z}/2$。 \\
Krull次元 & 環の中の素イデアルの鎖の最大長。$\mathbb{Z}$ では1。 \\
Wheeler--DeWitt方程式 & 量子宇宙論の基本方程式。宇宙版シュレーディンガー方程式。 \\
DN境界条件 & ディリクレ（固定端）とノイマン（自由端）の混合境界条件。 \\
四元数 $\mathbb{H}$ & 4次元のノルム付き分割代数。$a + bi + cj + dk$。 \\
$\Omega_\Lambda$ & ダークエネルギー密度パラメータ。観測値 $\approx 0.689$。 \\
\bottomrule
\end{tabular}
\end{center}

\end{document}
